\documentclass[twocolumn]{aastex631}
\begin{document}
\title{The reionization ionizing-photon-budget shortfall for star-forming galaxies is not robust to systematics at z\~{}7 (highC systematic corner)}
\author{NebulaMind Lab (autonomous pipeline)}
\affiliation{NebulaMind Lab, \url{https://lab.nebulamind.net}}
\begin{abstract}
Reionization ionizing-photon-budget reconciliation at z\~{}7: star-forming galaxies require f\_esc=0.154 (+0.144/-0.076) to close the budget (Madau-Dickinson SFRD, log xi\_ion=25.5+/-0.15, clumping C=2-5, JWST-SFRD tail), versus indirect-proxy-inferred f\_esc=0.062 (+0.108/-0.039) from LzLCS O32/beta calibrations. Median delta(required-inferred)=+0.079 dex-frac (16-84\%: -0.036 to +0.225); 77\% of the systematic MC shows a shortfall. Conclusion: the budget CLOSES within the systematic, and the sign holds under both O32 and beta calibrations. The result is bounded by the xi\_ion x clumping x proxy-calibration systematic, not by statistics. Generated autonomously from public data (jwst) via the NebulaMind Lab runner: a bounded, reproducible, descriptive study.
\end{abstract}
\section{Introduction}
Reconciling the ionizing photon budget during reionization has been a longstanding challenge in understanding the early universe. Recent studies have highlighted potential discrepancies between the observed star formation rate density (SFRD) and the required ionizing photons to drive reionization [Muñoz2024]. This has led to questions about whether star-forming galaxies alone can account for the necessary photons, or if additional sources are needed [Park2022]. To address this issue, we revisit the photon budget calculation using a literature-anchored approach.
\section{Data and method}
Our method relies on published values and calibrations from previous studies. We adopt the Madau \& Dickinson (2014) analytic fitting function for the cosmic SFRD, as well as the xi\_ion and O32/beta f\_esc proxy calibrations from LzLCS [Chisholm+22, Flury+22; Simmonds+24]. By combining these elements, we perform a systematic reconciliation of the ionizing photon budget at z\~{}7.
\section{Result}
Our calculation reveals that star-forming galaxies require an escape fraction of f\_esc=0.154 (+0.144/-0.076) to close the reionization photon budget, assuming a Madau-Dickinson SFRD, log xi\_ion=25.5+/-0.15, and clumping factor C=2-5. In contrast, indirect-proxy-inferred f\_esc from LzLCS O32/beta calibrations yields a value of 0.062 (+0.108/-0.039). The median difference between the required and inferred escape fractions is +0.079 dex-frac (16-84\%: -0.036 to +0.225), with 77\% of systematic Monte Carlo simulations showing a shortfall.
\begin{figure}\centering\includegraphics[width=\columnwidth]{result.png}\caption{The reionization ionizing-photon-budget shortfall for star-forming galaxies is not robust to systematics at z\~{}7 (highC systematic corner)}\end{figure}
\section{Caveats}
It is essential to acknowledge the limitations of our approach. Our calculation relies on automated, single-selection, uncalibrated measurements from published literature, which may introduce biases and uncertainties. The use of proxy calibrations for f\_esc and xi\_ion can also lead to discrepancies, as these values are subject to variation depending on the specific galaxy sample and observational conditions. Additionally, our analysis does not account for potential contributions from other ionizing sources, such as active galactic nuclei or quasars. Further research is needed to refine these estimates and better understand the complexities of reionization.
\section*{References}
\footnotesize
[Muoz2024] Muñoz et al. (2024). Reionization after JWST: a photon budget crisis?. bibcode 2024MNRAS.535L..37M \par [Park2022] Park et al. (2022). Calibrating excursion set reionization models to approximately conserve ionizing photons. bibcode 2022MNRAS.517..192P \par [Duncan2015] Duncan et al. (2015). Powering reionization: assessing the galaxy ionizing photon budget at z < 10. bibcode 2015MNRAS.451.2030D \par [Davies2021] Davies et al. (2021). The Predicament of Absorption-dominated Reionization: Increased Demands on Ionizing Sources. bibcode 2021ApJ...918L..35D \par [Madau2017] Madau (2017). Cosmic Reionization after Planck and before JWST: An Analytic Approach. bibcode 2017ApJ...851...50M

\end{document}
